% Unit 0: What are Differential Equations
% Condensed Cheat Sheet (Part I only)
% Prepared by Staples Education
% Content informed by the local curriculum modules (0.1 to 0.4) and the
% QUIZ_DATA micro-practice and unit-mastery items for Unit 0.

\documentclass[11pt]{article}

\usepackage[margin=1in]{geometry}
\usepackage{amsmath}
\usepackage{amssymb}
\usepackage{xcolor}
\usepackage[most]{tcolorbox}
\usepackage{enumitem}
\usepackage{hyperref}
\hypersetup{colorlinks=true, linkcolor=headerblue, urlcolor=headerblue}

\tcbuselibrary{skins}

% ---------------------------------------------------------------------------
% Unified style-guide palette (see LATEX_STYLE_GUIDE.md)
% ---------------------------------------------------------------------------
\definecolor{headerblue}{HTML}{1C569E}   % overview / structural framing
\definecolor{accentgreen}{HTML}{2E7D32}  % algorithms / methods
\definecolor{accentorange}{HTML}{E27A1E} % formulas / concept takeaways
\definecolor{workpurple}{HTML}{A020F0}   % worked-example tracking (vibrant)
\definecolor{warnred}{HTML}{C0392B}      % crimson pitfall / warning

% ---------------------------------------------------------------------------
% Subsection typography: headerblue numbered sub-module titles
% ---------------------------------------------------------------------------
\usepackage{titlesec}
\titleformat{\subsection}
  {\normalfont\large\bfseries\color{headerblue}}{\thesubsection}{1em}{}

% ---------------------------------------------------------------------------
% Six core custom boxes (unbreakable, title={#1} protected)
% ---------------------------------------------------------------------------

% 1. Blue overview capsule for the topics summary.
\newtcolorbox{topicsbox}{
    enhanced,
    colback=headerblue!6!white, colframe=headerblue,
    coltitle=white, fonttitle=\bfseries,
    title=Unit 0 at a Glance: Core Sub-Topics,
    boxrule=0.6pt, arc=2mm,
    left=3mm, right=3mm, top=2mm, bottom=2mm
}

% 2. Orange formula container for standard forms and master formulas.
\newtcolorbox{formulabox}[1][Master Formula]{
    enhanced,
    colback=accentorange!8!white, colframe=accentorange,
    coltitle=white, fonttitle=\bfseries,
    title={#1},
    boxrule=0.6pt, arc=1mm,
    left=3mm, right=3mm, top=2mm, bottom=2mm
}

% 3. Green algorithm block for explicitly laid-out procedures.
\newtcolorbox{methodbox}[1][Method]{
    enhanced,
    colback=accentgreen!6!white, colframe=accentgreen,
    coltitle=white, fonttitle=\bfseries,
    title={#1},
    boxrule=0.6pt, arc=1mm,
    left=3mm, right=3mm, top=2mm, bottom=2mm
}

% 4. Purple west-bordered worked example for step-by-step solutions.
\newtcolorbox{examplebox}[1][Worked Example]{
    enhanced,
    colback=workpurple!4!white, colframe=workpurple,
    coltitle=white, fonttitle=\bfseries,
    title={#1},
    boxrule=0.4pt, sharp corners,
    borderline west={3pt}{0pt}{workpurple},
    left=5mm, right=3mm, top=2mm, bottom=2mm
}

% 5. Orange thin-bordered final concept takeaway.
\newtcolorbox{conceptbox}[1][Key Takeaway]{
    enhanced,
    colback=accentorange!5!white, colframe=accentorange,
    coltitle=white, fonttitle=\bfseries,
    title={#1},
    boxrule=0.3pt, arc=1mm,
    left=3mm, right=3mm, top=2mm, bottom=2mm
}

% 6. Crimson west-bordered warning box for common pitfalls.
\newtcolorbox{warningbox}[1][Common Pitfall]{
    enhanced,
    colback=red!3!white, colframe=warnred,
    coltitle=white, fonttitle=\bfseries,
    title={#1},
    boxrule=0.4pt, sharp corners,
    borderline west={3pt}{0pt}{warnred},
    left=5mm, right=3mm, top=2mm, bottom=2mm
}

\setlist[enumerate]{itemsep=4pt, topsep=4pt}

% Number and index sections plus subsections (e.g., 0.1) two levels deep.
\setcounter{secnumdepth}{2}
\setcounter{tocdepth}{2}

\begin{document}

% ---------------------------------------------------------------------------
% Title block
% ---------------------------------------------------------------------------
\begin{center}
    {\LARGE \bfseries Unit 0: What are Differential Equations}\\[6pt]
    {\large Prepared by Staples Education}\\[4pt]
    {\itshape Condensed Cheat Sheet}
\end{center}

\vspace{0.4em}

% ---------------------------------------------------------------------------
% Topics overview box
% ---------------------------------------------------------------------------
\begin{topicsbox}
This opening unit builds the core mental model of the whole course: a
differential equation is an equation whose unknown is a \emph{function}, tied
to its own rates of change. Each sub-topic below sharpens that idea, moving
from raw definition toward why these equations describe the changing world.
\begin{itemize}[leftmargin=1.4em, itemsep=3pt]
    \item \textbf{What are Differential Equations} --- the unknown is a
          function, and the equation relates it to its own derivatives.
    \item \textbf{How Differential Equations Work} --- what it means for a
          function to be a solution, and why one equation yields a whole
          family of curves.
    \item \textbf{A Tourist's Guide} --- the vocabulary of order, linearity,
          ODE versus PDE, and the slope-field picture of a rate law.
    \item \textbf{Why You Are Learning Them} --- how laws about rates and
          forces become equations we can analyze, from growth to cooling.
\end{itemize}
\end{topicsbox}

\vspace{0.6em}
\hrule
\vspace{0.6em}

% ===========================================================================
% PART I: CONDENSED CHEAT SHEET
% ===========================================================================
\section*{Part I --- Condensed Cheat Sheet}
\setcounter{section}{0}   % subsections auto-number as the unit modules 0.1, 0.2, ...

\subsection{What are Differential Equations}

A \emph{differential equation} is an equation whose unknown is a function and
which relates that function to one or more of its own derivatives. This is the
one structural fact that separates it from algebra --- in algebra you solve for
a number, but here you solve for a whole function. The \emph{order} of the
equation is simply the highest derivative that appears.

\begin{formulabox}[Definition and Order]
A differential equation links an unknown function $y(x)$ to its derivatives:
\[
    F\!\left(x,\; y,\; \frac{dy}{dx},\; \frac{d^{2}y}{dx^{2}},\; \dots\right) = 0.
\]
The \textbf{order} is the highest derivative present. For example,
$\dfrac{dy}{dx} = 3y$ is \emph{first order}, while $\dfrac{d^{2}y}{dx^{2}} + 4y = 0$
is \emph{second order}.
\end{formulabox}

\begin{warningbox}[Common Pitfall: The Unknown Is a Function, Not a Number]
The most common first-week mistake is hunting for a single number, as in
algebra. The unknown here is an entire function $y(x)$ --- a rule defined across
its whole domain. A solution is satisfied at \emph{every} point, not at one
isolated value.
\end{warningbox}

\begin{examplebox}[Worked Example 0.1: Read $\frac{dy}{dx} = 3y$]
\textbf{Unknown.} The function $y(x)$, not the constant $3$ and not the input $x$.\\
\textbf{Order.} The highest derivative is the first derivative, so the equation
is first order.\\
\textbf{Reading.} It says: find every function whose rate of change equals three
times its own current value.
\end{examplebox}

\subsection{How Differential Equations Work}

To \emph{solve} a differential equation is to find the function --- or family of
functions --- that satisfies it for all inputs. You \emph{check} a candidate by
substituting it and its derivatives back in and confirming both sides agree
everywhere. Because a first-order equation fixes only the slope at each point
and not the starting height, its general solution carries one arbitrary
constant, and an initial condition selects one curve from that family.

\begin{formulabox}[Solution, General versus Particular]
A function is a solution when substitution makes the equation hold for all $x$.
The \textbf{general solution} of a first-order equation carries one free
constant; an \textbf{initial condition} $y(x_0) = y_0$ pins it down:
\[
    y = F(x) + C
    \;\;\xrightarrow{\;y(x_0)=y_0\;}\;\;
    \text{one particular curve.}
\]
\end{formulabox}

\begin{methodbox}[Algorithm: Verifying a Candidate Solution]
\begin{enumerate}[leftmargin=1.6em, itemsep=2pt]
    \item Compute the derivatives of the candidate that the equation requires.
    \item Substitute the candidate and those derivatives into the equation.
    \item Simplify the left-hand side and the right-hand side independently.
    \item Confirm the two sides are equal as an identity in $x$ --- true for
          every input, not just one point.
\end{enumerate}
\end{methodbox}

\begin{examplebox}[Worked Example 0.2: Verify $y = e^{2x}$ solves $\frac{dy}{dx} = 2y$]
\textbf{Step 1.} Differentiate the candidate:
\[ \frac{dy}{dx} = 2e^{2x}. \]
\textbf{Step 2.} Compute the right side from the candidate:
\[ 2y = 2e^{2x}. \]
\textbf{Step 3.} Both sides equal $2e^{2x}$ for all $x$, so $y = e^{2x}$ is a
solution.
\end{examplebox}

\subsection{A Tourist's Guide to Differential Equations}

Before solving anything, you learn to classify. An equation is an \emph{ordinary}
differential equation (ODE) when its unknown depends on a single independent
variable, and \emph{partial} (PDE) when it depends on several. It is
\emph{linear} when the unknown and its derivatives appear only to the first
power and are never multiplied together. Geometrically, a first-order equation
is a \emph{slope field} --- a direction mark at every point that solution curves
must follow.

\begin{formulabox}[Classification Vocabulary]
\begin{itemize}[leftmargin=1.4em, itemsep=2pt]
    \item \textbf{Order} --- the highest derivative present.
    \item \textbf{ODE vs PDE} --- one independent variable versus several.
    \item \textbf{Linear} --- $y$ and its derivatives appear to the first power
          only, with no products such as $y\,y'$ and no terms like $\sin y$.
    \item \textbf{Slope field} --- the picture $\dfrac{dy}{dx} = f(x,y)$ paints,
          assigning a direction to every point of the plane.
\end{itemize}
\end{formulabox}

\begin{warningbox}[Common Pitfall: Order Is Not Term Count or Degree]
The order is set \emph{only} by the highest derivative --- not by the number of
terms, the size of the coefficients, or the power to which the function is
raised. The equation $\left(\dfrac{dy}{dx}\right)^{3} = y$ is still
\emph{first order}; the cube is a degree, not an order.
\end{warningbox}

\begin{examplebox}[Worked Example 0.3: Classify three equations]
\textbf{(a)} $\dfrac{dy}{dx} + y = x$ --- first order, linear, ODE.\\
\textbf{(b)} $\dfrac{d^{2}\theta}{dt^{2}} + \sin\theta = 0$ --- second order,
\emph{nonlinear} (the $\sin\theta$ term), ODE.\\
\textbf{(c)} $\dfrac{\partial u}{\partial t} = \dfrac{\partial^{2} u}{\partial x^{2}}$
--- a PDE, since $u$ depends on both $t$ and $x$.
\end{examplebox}

\subsection{Why You Are Learning Differential Equations}

Differential equations are the language of science because most physical laws
describe \emph{how} a quantity changes, which is exactly what a derivative
measures. To model a phenomenon you identify the changing quantity and what its
rate of change depends on. The same handful of rate laws --- a rate proportional
to an amount, or to a difference --- reappear across growth, decay, and cooling.

\begin{formulabox}[Standard Models]
A rate proportional to the current amount gives exponential behavior:
\[
    \frac{dP}{dt} = kP \;\Longrightarrow\; P(t) = P_{0}e^{kt}
    \qquad (k>0 \text{ grows},\; k<0 \text{ decays}).
\]
Radioactive decay and Newton's law of cooling share the structure:
\[
    \frac{dN}{dt} = -\lambda N,
    \qquad
    \frac{dT}{dt} = -k\,(T - T_{s}).
\]
Newton's second law writes acceleration as a second derivative of position:
$F = m\dfrac{d^{2}x}{dt^{2}}$.
\end{formulabox}

\begin{examplebox}[Worked Example 0.4: Set up a growth model]
\textbf{Step 1.} A population whose growth rate is proportional to its size
obeys
\[ \frac{dP}{dt} = kP. \]
\textbf{Step 2.} Check that $P(t) = P_{0}e^{kt}$ fits: its derivative is
$kP_{0}e^{kt} = kP$, matching the equation, and $P(0) = P_{0}$ sets the initial
size. The single rate law captures the whole exponential family.
\end{examplebox}

\vspace{0.8em}
\begin{conceptbox}[Unit 0 Key Takeaway]
A differential equation trades the algebra question ``what number?'' for the
calculus question ``what \textcolor{accentorange}{function}?'' It pins down the
\textcolor{accentorange}{rate of change} at every point, so its general solution
is a \textcolor{accentorange}{whole family of curves} that an
\textcolor{accentorange}{initial condition} narrows to one. Read the structure
first --- order, linearity, and what is changing --- and every later technique
has a place to land.
\end{conceptbox}

\end{document}
